%! Composition of Functions — Unit 1, Lesson 1.5 — Exit Ticket (Answer Key)
\documentclass[10pt]{article}
\usepackage{precalculus-article}
\usepackage{precalculus-key}   % pulls in -boxes; do NOT also load -boxes
\newcommand{\TallMath}[1]{$\displaystyle #1\rule[-1.4em]{0pt}{3.2em}$}

\begin{document}

\pageheader{Unit 1, Lesson 1.5}{Exit Ticket --- Answer Key}

\vspace{0.12in}

\textit{Let $f(x) = x + 4$ and $g(x) = 3x$ for items 1 and 2.}

\vspace{0.1in}

\begin{enumerate}[itemsep=12pt]

\item Evaluate both composites at $x = 2$. Write the middle number down.

\begin{work}
  f(g(2)) &= f(6) \\
          &= 6 + 4 \\
          &= 10 \\
  g(f(2)) &= g(6) \\
          &= 3(6) \\
          &= 18
\end{work}

$f(g(2)) = $ \ans{10} \qquad $g(f(2)) = $ \ans{18}

\item Write each composite as a simplified formula.

\begin{work}
  (f \circ g)(x) &= f(3x) \\
                 &= 3x + 4 \\
  (g \circ f)(x) &= g(x + 4) \\
                 &= 3x + 12
\end{work}

$(f \circ g)(x) = $ \ans{$3x + 4$} \qquad $(g \circ f)(x) = $ \ans{$3x + 12$}

\smallskip
Are they the same function? \ans{No}

\item Decompose $h(x) = (x - 5)^3$ into an inner and an outer function.

\smallskip
Inner: $g(x) = $ \ans{$x - 5$} \qquad Outer: $f(u) = $ \ans{$u^3$}

\end{enumerate}

\end{document}
